% Generated draft: requires mathematical and source review.
\documentclass[11pt]{article}
\usepackage[margin=25mm]{geometry}
\usepackage{fontspec}
\setmainfont{Latin Modern Roman}
\usepackage{amsmath,amssymb,mathtools,unicode-math}
\setmathfont{Latin Modern Math}
\usepackage{enumitem,needspace,xurl,hyperref}
\setlength{\parindent}{0pt}
\setlength{\parskip}{5pt}
\setlength{\emergencystretch}{3em}
\newcommand{\N}{\mathbb{N}}
\newcommand{\Z}{\mathbb{Z}}
\newcommand{\Q}{\mathbb{Q}}
\newcommand{\R}{\mathbb{R}}
\newcommand{\C}{\mathbb{C}}
\newcommand{\F}{\mathbb{F}}
\newcommand{\A}{\mathbb{A}}
\newcommand{\PP}{\mathbb{P}}
\newcommand{\E}{\mathbb{E}}
\newcommand{\eps}{\varepsilon}
\newcommand{\dd}{\,\mathrm{d}}
\begin{document}
\section*{UnsolvedMath ID: 701668}
\textbf{Source problem code:} catalogue entry
\subsection{Definitions and mathematical statement}
Let \(\N\) be the countably infinite index set. A nonprincipal ultrafilter \(\mathscr U\) on \(\N\) is a family of subsets of \(\N\) that contains \(\N\), excludes \(\varnothing\), is closed under finite intersections and passage to supersets, contains exactly one of \(A\) and \(\N\setminus A\) for every \(A\subseteq\N\), and contains no finite set.

For any sequence \((G_n)_{n\in\N}\) of groups, define
\[
N_{\mathscr U}=\left\{(g_n)\in\prod_{n\in\N}G_n:\{n\in\N:g_n=1_{G_n}\}\in\mathscr U\right\}.
\]
The ultraproduct with respect to \(\mathscr U\) is the quotient group
\[
\prod_{n\in\N}G_n/\mathscr U
:=\left(\prod_{n\in\N}G_n\right)/N_{\mathscr U}.
\]
Let \(\mathcal H(\mathscr U)\) denote the class of all groups \(H\) for which there exist a sequence \((G_n)_{n\in\N}\) of groups and a surjective group homomorphism
\[
\prod_{n\in\N}G_n/\mathscr U\twoheadrightarrow H.
\]
There are no finiteness or cardinality restrictions on the groups \(G_n\) or \(H\). The representing sequence may depend on \(H\) and on the ultrafilter.

\textbf{(a)} Is the following assertion true?
\[
\text{For every pair of nonprincipal ultrafilters }\mathscr U,\mathscr U'\text{ on }\N,
\qquad \mathcal H(\mathscr U)\subseteq\mathcal H(\mathscr U').
\]
Since the pair is arbitrary, this assertion is equivalent to equality of these classes for every such pair.

\textbf{(b)} If the answer to (a) is negative, the source asks whether the following assertion nevertheless holds:
\[
\begin{gathered}
\text{For every pair of nonprincipal ultrafilters }\mathscr U,\mathscr U'\text{ on }\N,\\
\text{there exists a nonprincipal ultrafilter }\mathscr U''\text{ on }\N\text{ such that}\\
\mathcal H(\mathscr U)\cup\mathcal H(\mathscr U')
\subseteq\mathcal H(\mathscr U'').
\end{gathered}
\]
Here \(\mathscr U''\) must work simultaneously for every group in the displayed union; it may depend on \(\mathscr U\) and \(\mathscr U'\), but not on the individual group.

The source states that an affirmative answer to (b) would imply that the class
\[
\mathcal C=\bigcup_{\mathscr U\text{ nonprincipal on }\N}\mathcal H(\mathscr U)
\]
is closed under finite direct products: for every positive integer \(r\) and all \(H_1,\ldots,H_r\in\mathcal C\), one has \(\prod_{i=1}^r H_i\in\mathcal C\). Here ``countable ultraproduct'' refers to the index set, not to the cardinalities of its factors.

\textbf{Supplied status qualification.} The registry labels the problem partially solved. Its discussion labels (b) solved, but the result actually described there is that an affirmative answer to (b) is consistent with the ZFC axioms, attributed to S. M. Corson (2025). This supplied consistency report is not an assertion here that (b) has been proved in ZFC or that (a) has been resolved. The exact hypotheses and scope of the reported consistency result require source review.
\subsection{Short English statement}
For each nonprincipal ultrafilter on \(\N\), consider all groups obtainable as quotients of ultraproducts of arbitrary groups indexed by \(\N\). Part (a) asks whether this class is independent of the ultrafilter. If it is not, part (b) asks whether any two such classes are contained in a third class arising from one nonprincipal ultrafilter. The supplied update reports consistency with ZFC of an affirmative answer to (b), rather than an unconditional ZFC proof.
\subsection{Sources}
\begin{itemize}
\item Local UnsolvedMath registry, permanent ID \(701668\), title \emph{Kourovka Notebook Problem 20.10}, exact source code \texttt{catalogue entry}; supplied record with upstream ID \(701668\) and revision \texttt{d371ee917ac9647fbb06cd6fc529deec6f3dda82}. The registry supplies the status \texttt{partially\_solved}.
\item G. M. Bergman, proposer, 2022. Source attribution: \emph{The Kourovka Notebook (Unsolved Problems in Group Theory)}, 20th issue (2022), problem 20.10; also the supplied 21st issue edition (2026), arXiv:\texttt{1401.0300v48}. The supplied collection description names E. I. Khukhro and V. D. Mazurov as editors.
\item G. M. Bergman, \emph{Pacific J. Math.} \textbf{274} (2015), 451--495. The statement cites this reference in connection with the implication from an affirmative answer to (b) to closure under finite direct products. No article title or link is supplied.
\item S. M. Corson, \emph{Preprint}, 2025, \url{https://arxiv.org/pdf/2503.09228}. The supplied discussion reports that an affirmative answer to (b) is consistent with the ZFC axioms of set theory and attaches the label ``[Solved]'' to (b). No preprint title is supplied.
\item Supplied dated curation: the status review is dated \texttt{2026-10-16}, identifies its evidence as \texttt{maintained\_source}, and attributes its report to the 21st issue of the Kourovka Notebook (October 2026, arXiv:\texttt{1401.0300v48}). Its notes describe the subproblem as ``confirmed solved (*)'' while specifying the consistency result above. This is supplied dated curation, not an independently verified literature assessment; no source visit is claimed.
\end{itemize}
\end{document}
